\documentclass[11pt,a4paper]{article}

% ---------------------------------------------------------
% PREAMBLE (stable, warning-free, Overleaf-safe)
% ---------------------------------------------------------
\usepackage[utf8]{inputenc}
\usepackage[T1]{fontenc}
\usepackage{lmodern}
\usepackage{amsmath, amssymb, amsfonts}
\usepackage{bm}
\usepackage{geometry}
\usepackage{graphicx}
\usepackage{hyperref}
\usepackage{cite}
\usepackage{authblk}
\usepackage{physics}
\usepackage{mathtools}

\geometry{margin=2.4cm}

\title{\textbf{Companion XLV: Weak Lensing in the Relaxation-Driven Cyclic Cosmology}}
\author{Michael Lehmann \\ \texttt{mi.lehmann@gmx.de}}
\date{April 2026}

\begin{document}
\maketitle

\begin{abstract}
Weak gravitational lensing probes the integrated gravitational potential along the 
line of sight and provides a direct measurement of the geometry and growth of 
structure in the Universe. In the standard cosmological model, the lensing field is 
determined by the sum of the two Newtonian potentials, which evolve identically in 
the absence of anisotropic stress. In the Relaxation-Driven Cyclic Cosmology (RDCC), 
the scale-dependent evolution of the gravitational potential modifies the lensing 
kernel, the shear field, and the convergence power spectrum. This Companion develops 
a continuous description of weak lensing in RDCC, deriving how the relaxation scale 
influences the deflection of light, the geometry of shear patterns, and the 
large-scale coherence of the lensing field. The resulting framework provides a 
distinctive observational signature of the RDCC gravitational field.
\end{abstract}
% ---------------------------------------------------------
\section{Introduction}

Weak gravitational lensing arises from the deflection of light by the gravitational 
potential of large-scale structure. The resulting distortions of galaxy shapes 
encode the integrated gravitational field along the line of sight. In the standard 
cosmological model, the lensing potential is given by the sum of the two Newtonian 
potentials, which evolve identically in the absence of anisotropic stress.

In RDCC, the gravitational potential evolves differently. The relaxation mechanism 
suppresses the decay of small-scale modes and enhances the coherence of large-scale 
modes. This modifies the lensing potential, the shear field, and the convergence 
power spectrum in a scale-dependent manner. Weak lensing becomes a direct probe of 
the relaxation scale $k_{\mathrm{rel}}$ and the temporal structure of the RDCC 
gravitational field.
% ---------------------------------------------------------
\section{The RDCC Lensing Potential}

The lensing potential is defined as
\begin{equation}
    \phi(\hat{\mathbf{n}})
    = -2 \int_{0}^{\chi_{s}}
      d\chi\,
      \frac{\chi_{s} - \chi}{\chi_{s}\chi}\,
      \left[\Phi + \Psi\right](\chi,\hat{\mathbf{n}}).
\end{equation}

In RDCC, the metric potentials evolve differently:
\begin{align}
    \Phi_{\mathrm{RDCC}}(k,a)
    &= \Phi(k,a)
       \left[
         1 - \chi_{\Phi}
         \left(\frac{k}{k_{\mathrm{rel}}}\right)^{\alpha}
       \right], \\
    \Psi_{\mathrm{RDCC}}(k,a)
    &= \Psi(k,a)
       \left[
         1 - \chi_{\Psi}
         \left(\frac{k}{k_{\mathrm{rel}}}\right)^{\alpha}
       \right].
\end{align}

The lensing potential becomes
\begin{equation}
    \phi_{\mathrm{RDCC}}
    = \phi
      \left[
        1 - \chi_{\phi}
        \left(\frac{k}{k_{\mathrm{rel}}}\right)^{\alpha}
      \right].
\end{equation}

This modification suppresses small-scale lensing power and enhances large-scale 
coherence.
% ---------------------------------------------------------
\section{Shear Geometry and the RDCC Shear Field}

The shear field is defined as the second angular derivative of the lensing 
potential:
\begin{equation}
    \gamma = \frac{1}{2}
             \left(
               \partial_{1}^{2}
               - \partial_{2}^{2}
             \right)\phi
             + i\,\partial_{1}\partial_{2}\phi.
\end{equation}

In RDCC, the modified lensing potential produces a shear field with smoother 
small-scale structure and enhanced large-scale coherence. The shear correlation 
function becomes
\begin{equation}
    \xi_{\pm,\mathrm{RDCC}}(\theta)
    = \xi_{\pm}(\theta)
      \left[
        1 - \chi_{\gamma}
        \left(\frac{1}{k_{\mathrm{rel}}\theta}\right)^{\alpha}
      \right].
\end{equation}

This modification produces a characteristic RDCC signature in cosmic shear surveys.
% ---------------------------------------------------------
\section{Convergence Power Spectrum}

The convergence field is defined as
\begin{equation}
    \kappa
    = \frac{1}{2}\nabla^{2}\phi.
\end{equation}

The convergence power spectrum is given by
\begin{equation}
    C_{\ell}^{\kappa\kappa}
    = \int d\chi\,
      \frac{W^{2}(\chi)}{\chi^{2}}\,
      P_{\Phi+\Psi}\left(k=\frac{\ell}{\chi},\chi\right).
\end{equation}

In RDCC, the potential power spectrum becomes
\begin{equation}
    P_{\Phi+\Psi,\mathrm{RDCC}}(k)
    = P_{\Phi+\Psi}(k)
      \left[
        1 - \chi_{\kappa}
        \left(\frac{k}{k_{\mathrm{rel}}}\right)^{\alpha}
      \right].
\end{equation}

This modification suppresses small-scale convergence power and enhances large-scale 
structure, producing a distinctive RDCC signature.
% ---------------------------------------------------------
\section{Weak Lensing Tomography and the RDCC Kernel}

Weak lensing tomography probes the redshift evolution of the lensing field. In 
RDCC, the lensing kernel becomes
\begin{equation}
    W_{\mathrm{RDCC}}(\chi)
    = W(\chi)
      \left[
        1 - \chi_{W}
        \left(\frac{k(\chi)}{k_{\mathrm{rel}}}\right)^{\alpha}
      \right].
\end{equation}

This modification alters the redshift dependence of the lensing signal, producing 
a smoother evolution and enhanced coherence across tomographic bins.

These signatures provide a direct observational probe of the relaxation scale.
% ---------------------------------------------------------
\section{Conclusion}

Weak gravitational lensing in the Relaxation-Driven Cyclic Cosmology reflects the 
scale-dependent evolution of the gravitational potential. The relaxation mechanism 
modifies the lensing potential, the shear field, and the convergence power 
spectrum, producing distinctive signatures in cosmic shear surveys. These effects 
provide a direct observational window into the relaxation scale and the temporal 
structure of the RDCC gravitational field. The present Companion establishes the 
theoretical foundation for future studies of lensing tomography, shear statistics, 
and the interplay between matter and light in RDCC.

% ========================
% References
% ========================

\begin{thebibliography}{10}

\bibitem{Flagship} Lehmann, M. (2026). Relaxation-Driven Cyclic Cosmology (RDCC): A Minimal CPT-Symmetric Framework with a Single Parameter. Flagship Paper. Zenodo: \url{https://zenodo.org/records/19744958} (v43.0 or later).

\bibitem{Zenodo} The complete RDCC companion series (I--LVII) is available at the same Zenodo record above.

% Thematisch relevante Companions
\bibitem{CompXVI} Companion XVI: Boltzmann Code and Modified Hierarchy.
\bibitem{CompXXXI} Companion XXXI: RDCC Halo Model and Non-Linear Structure.
\bibitem{CompXXXII} Companion XXXII--XXXIX: Galaxy--Halo Connection, Cosmic Web, Lensing, RSD, ISW, tSZ, Rotation Curves, CGM.

% Wichtige externe Referenzen
\bibitem{Planck2020} Planck Collaboration (2020), Astron. Astrophys. 641, A6.
\bibitem{DESI2024} DESI Collaboration (2024/2025), arXiv: [aktuelle $f\sigma_8$/BAO-Paper].
\bibitem{JWST2024} Maiolino et al. (2024), Nature (JWST high-$z$ galaxies/quasars).
\bibitem{Kaiser1987} Kaiser, N. (1987), Mon. Not. R. Astron. Soc. 227, 1 (RSD).
\bibitem{Bartelmann2001} Bartelmann, M. \& Schneider, P. (2001), Phys. Rep. 340, 291 (Weak Lensing Review).
\bibitem{HuOkamoto2002} Hu, W. \& Okamoto, T. (2002), Astrophys. J. 574, 566 (CMB Lensing).
\bibitem{LewisChallinor2006} Lewis, A. \& Challinor, A. (2006), Phys. Rep. 429, 1 (CMB Lensing Review).
\bibitem{NFW1997} Navarro, J. F., Frenk, C. S. \& White, S. D. M. (1997), Astrophys. J. 490, 493 (Halo Profiles).

\end{thebibliography}

\end{document}
